\clearpage
\SourceLeaf{169}%
\pdfbookmark[0]{Göttingische gelehrte Anzeigen. 1817 März 10.}{notice1817march10}
\begin{center}
\rule{.82\linewidth}{.35pt}\par
Göttingische gelehrte Anzeigen.\quad 1817 März 10.\par
\rule{.82\linewidth}{.35pt}
\end{center}
Am 10. Februar reichte Herr Hofrat \textsc{Gauss} der Königlichen Gesellschaft eine Vorlesung mit dem Titel ein:
\begin{center}\itshape
Theorematis fundamentalis in doctrina de residuis quadraticis demonstrationes et ampliationes novae.
\end{center}

Es ist eine Besonderheit der höheren Arithmetik, dass sich so viele ihrer schönsten Sätze mit größter Leichtigkeit durch die Untersuchung von Zahlenbeispielen entdecken lassen, während ihre Beweise keineswegs naheliegen. Oft findet man diese erst nach vielen vergeblichen Versuchen mithilfe tiefgehender Untersuchungen und glücklicher Kombinationen. Dieses bemerkenswerte Phänomen entspringt der oft wunderbaren Verknüpfung der verschiedenartigen Lehren in diesem Teil der Mathematik. Eben deshalb lassen sich Sätze, für die man anfangs jahrelang vergeblich einen Beweis gesucht hatte, später häufig auf mehreren ganz verschiedenen Wegen beweisen. Sobald ein neuer Satz anhand von Zahlenbeispielen entdeckt ist, muss man die Auffindung \emph{irgendeines} Beweises freilich als erstes Erfordernis ansehen. Ist dies aber gelungen, so darf man in der höheren Arithmetik die Untersuchung nicht immer für abgeschlossen und die Suche nach anderen Beweisen für überflüssigen Luxus halten. Denn zum einen findet man gewöhnlich die schönsten und einfachsten
\SourceLeaf{170}%
Beweise nicht zuerst; zum anderen macht gerade die Einsicht in die wunderbare Verknüpfung der Wahrheit der höheren Arithmetik einen Hauptreiz dieses Studiums aus und führt nicht selten wiederum zur Entdeckung neuer Wahrheiten. Aus diesen Gründen ist hier die Auffindung neuer Beweise für schon bekannte Wahrheiten oft als mindestens ebenso wichtig anzusehen wie die Entdeckung der Wahrheit selbst. Kennern der höheren Arithmetik sind diese Überlegungen nicht neu. Bekanntlich besteht ein großer Teil von \textsc{Eulers} Verdiensten um sie darin, Beweise für Sätze gefunden zu haben, die schon \textsc{Fermat}, wie es scheint anhand von Zahlenbeispielen, entdeckt hatte.

Die Lehre von den quadratischen Resten veranschaulicht das eben Gesagte. Sie beruht hauptsächlich auf dem sogenannten Fundamentaltheorem: Die wechselseitigen Beziehungen zweier (ungerader positiver) Primzahlen zueinander (insofern die eine quadratischer Rest oder Nichtrest der anderen ist) sind gleich, wenn eine oder beide Primzahlen die Form $4k+1$ haben, dagegen entgegengesetzt, wenn beide Primzahlen die Form $4k+3$ haben. Für Leser, die mit der höheren Arithmetik weniger vertraut sind, erinnern wir daran, dass eine ganze Zahl quadratischer Rest einer anderen heißt, wenn die erste durch Addition eines Vielfachen der zweiten ein Quadrat ergeben kann; Nichtrest dagegen, wenn dies unmöglich ist. Die Geschichte dieses schönen Satzes, der sich anhand von Zahlenbeispielen äußerst leicht finden lässt, wollen wir hier nicht vollständig wiederholen. Wir bemerken nur, dass der Verfasser der vorliegenden Abhandlung nach anfangs ziemlich lange vergeblich angestellten Untersuchungen nach und nach bereits vier ganz verschiedene Beweise gegeben hat. Zwei davon stehen in den \emph{Disquisitionibus Arithmeticis}; der dritte bildet den Gegenstand einer eigenen Abhandlung im sechzehnten Band der Commentationes, und der vierte ist in die Abhandlung \emph{summatio quarumdam serierum singularium} im ersten Band der \emph{Commentationes recentiores} eingearbeitet. Zu diesen beiden Abhandlungen kann man unsere Besprechungen von 1808, Mai 12 und September 19, heranziehen; dort finden sich auch ausführlichere historische Nachweise. Dass der Verfasser bei diesen vier Beweisen nicht stehen geblieben ist, obwohl jeder für sich an Strenge nichts zu wünschen übrig lässt, bedarf bei den Freunden der höheren Arithmetik keiner Rechtfertigung. Dennoch hätte er sich wahrscheinlich nicht so eifrig bemüht, noch weitere Beweise hinzuzufügen, wenn
\SourceLeaf{171}%
ihn nicht ein besonderer Umstand dazu veranlasst hätte, der hier erwähnt werden muss. Seit dem Jahr 1805 hatte er begonnen, sich mit den Theorien der kubischen und biquadratischen Reste zu beschäftigen, die noch weitaus reichhaltiger und interessanter sind als die Theorie der quadratischen Reste. Bei diesen Untersuchungen zeigten sich dieselben Erscheinungen wie bei jener, nur gleichsam in vergrößertem Maßstab. Sobald der richtige Weg eingeschlagen war, fand man anhand von Zahlenbeispielen sogleich eine Anzahl höchst einfacher Sätze, die diese Theorien ganz ausschöpfen, den Sätzen über quadratische Reste überraschend ähnlich sind und insbesondere auch ein Gegenstück zum Fundamentaltheorem bieten. Doch waren hier die Schwierigkeiten, völlig befriedigende Beweise für diese Sätze zu finden, noch viel größer. Erst nach vielen Versuchen, die über eine beträchtliche Reihe von Jahren fortgesetzt wurden, erreichte der Verfasser schließlich sein Ziel. Die große Analogie zwischen den Sätzen über quadratische und über höhere Reste ließ vermuten, dass es auch für beide analoge Beweise geben müsse. Die zuerst für quadratische Reste gefundenen Beweisarten ließen sich jedoch auf höhere Reste überhaupt nicht anwenden; gerade dies war der Beweggrund, für die ersteren immer weitere neue Beweise zu suchen. Der Verfasser wünscht daher, dass man die vorliegende Abhandlung, die noch einige neue Hilfsmittel für die Theorie der quadratischen Reste erschließt, als Vorläuferin der Theorie der kubischen und biquadratischen Reste betrachtet. Diese will er künftig veröffentlichen; ihre Gegenstände gehören zu den schwierigsten der höheren Arithmetik.

Die vorliegende Abhandlung besteht aus drei voneinander unabhängigen Teilen. Sie enthält den fünften und sechsten Beweis des Fundamentaltheorems und eine neue, mit dem dritten Beweis zusammenhängende Methode,\textperiodcentered\ um zu entscheiden, ob eine gegebene ganze Zahl quadratischer Rest oder Nichtrest einer gegebenen Primzahl ist. Unter den ersten vier Beweisen war der dritte unstreitig derjenige, der die größte Einfachheit mit der Unabhängigkeit von andersartigen Untersuchungen verband. Deshalb hat ihn auch \textsc{Legendre} in die neue Ausgabe seines \emph{Essai d'une théorie des nombres} aufgenommen. Der \emph{fünfte} Beweis scheint dem dritten in beiden Hinsichten mindestens gleichzukommen. Beide Beweise sind insofern verwandt, als sie von demselben Satz ausgehen; in ihrer weiteren Durchführung sind sie jedoch völlig von\SourceLeaf{172}%
einander verschieden. Dieser Satz lautet: Es sei $m$ eine (positive ungerade) Primzahl und $M$ eine ganze, durch $m$ nicht teilbare Zahl. Wenn ferner unter den Resten, die bei Division der Produkte
\[
M,\ 2M,\ 3M,\ 4M\mathinner{.\,.\,.\,.\,.\,.}\tfrac12(m-1)M
\]
durch $m$ entstehen, die Anzahl derjenigen, die größer als $\frac12m$ sind, mit $n$ bezeichnet wird, so ist $M$ quadratischer Rest oder Nichtrest von $m$, je nachdem $n$ gerade oder ungerade ist. Um nun zum Beweis des Fundamentaltheorems zu gelangen, nimmt man an, dass auch $M$ eine ungerade positive Primzahl ist und $N$ bezüglich $M$ und $m$ dasselbe bedeutet wie $n$ bezüglich $m$ und $M$. Ob $N$ gerade oder ungerade ist, entscheidet also, ob $m$ quadratischer Rest oder Nichtrest von $M$ ist. Durch eine sehr kurze Folge von Schlüssen zeigt der Verfasser, dass die Anzahl aller positiven ganzen Zahlen, die zugleich kleiner als $\frac12mM$ sind, bei Division durch $m$ einen Rest kleiner als $\frac12m$ und bei Division durch $M$ einen Rest kleiner als $\frac12M$ lassen,
\[
=\tfrac18(m-1)(M-1)+\tfrac12n+\tfrac12N
\]
und folglich stets
\[
\tfrac14(m-1)(M-1)+n+N
\]
eine gerade Zahl ist. Hat also mindestens eine der Zahlen $m,M$ die Form $4k+1$, sodass $\frac14(m-1)(M-1)$ gerade ist, so ist auch $n+N$ gerade; folglich sind $n$ und $N$ entweder beide gerade oder beide ungerade. Haben dagegen sowohl $m$ als auch $M$ die Form $4k+3$, so ist $n+N$ notwendig ungerade, also eine der Zahlen $n,N$ gerade und die andere ungerade. Zusammen mit dem obigen Satz folgt daraus unmittelbar das Fundamentaltheorem.

Der \emph{sechste} Beweis ist ebenso kurz und elegant wie der fünfte, beruht aber auf etwas kunstvolleren Kombinationen. Der beschränkte Raum dieser Blätter erlaubt nur, unter Weglassung der Einzelheiten, den Hauptgedanken zu berühren. Es bezeichnen:
\begin{description}
\item[$p,q$] zwei (verschiedene positive ungerade) Primzahlen,
\item[$\alpha$] eine sogenannte \emph{primitive Wurzel} zum Modul $p$, also eine durch $p$ nicht teilbare (hier positive) ganze Zahl, für die keine niedrigere Potenz als $\alpha^{p-1}$ modulo $p$ zu eins kongruent ist
\item[$x$] eine unbestimmte Größe
\SourceLeaf{173}%
\item[$\xi$] die Funktion
\[
x-x^\alpha+x^\zeta-x^\eta+x^\theta-\text{usw.}-x^\lambda
\]
wobei der bequemeren Drucklegung wegen $\zeta,\eta,\theta\ldots\lambda$ anstelle der Zahlen $\alpha\alpha,\allowbreak \alpha^3,\allowbreak \alpha^4\ldots\allowbreak \alpha^{p-2}$ stehen;
\item[$\varepsilon$] die Einheit, positiv genommen, wenn $p$ die Form $4k+1$, negativ, wenn $p$ die Form $4k+3$ hat;
\item[$\delta$] die Einheit, positiv genommen, wenn mindestens eine der Zahlen $p,q$ die Form $4k+1$ hat, negativ, wenn beide die Form $4k+3$ haben;
\item[$\gamma$] die Einheit, positiv genommen, wenn $q$ quadratischer Rest von $p$ ist, negativ, wenn $q$ quadratischer Nichtrest von $p$ ist;
\item[$\beta$] die Einheit, positiv genommen, wenn $p$ quadratischer Rest von $q$, negativ, wenn $p$ quadratischer Nichtrest von $q$ ist.
\end{description}

Nach diesen Vorbereitungen folgt leicht aus Art. 51 der \emph{Disquisitiones Arithmeticae}, dass die Funktion
\[
\xi^q-x^q+x^{q\alpha}-x^{q\zeta}+x^{q\eta}-x^{q\theta}+\text{usw.}+x^{q\lambda}
\]
nach dem Ausmultiplizieren nur durch $q$ teilbare Koeffizienten hat. Setzt man diese Funktion $=qX$, so ist $X$ ein \emph{Polynom}, dessen Koeffizienten ebenfalls ganze Zahlen sind. Durch Schlüsse, auf die hier näher einzugehen zu weit führen würde, wird in der Abhandlung bewiesen, dass die Funktion $qX\xi$ bei Division durch $x^{p-1}+x^{p-2}+x^{p-3}+x^{p-4}+\text{usw.}+x+1$ den Rest
\[
\varepsilon p(\delta p^{\frac12(q-1)}-\gamma)
\]
lässt, dass also bei Division der Funktion $X\xi$ durch denselben Divisor der Rest
\[
\frac{\varepsilon p(\delta p^{\frac12(q-1)}-\gamma)}{q}
\]
entsteht. Diese Größe muss daher eine ganze Zahl sein. Da $\delta\delta=1$ ist, folgt daraus leicht, dass
\[
p^{\frac12(q-1)}-\gamma\delta
\]
durch $q$ teilbar sein muss. Nach einem bekannten Satz ist auch $p^{\frac12(q-1)}-\beta$ durch $q$ teilbar. Somit muss $\beta=\gamma\delta$ gelten, woraus wiederum unmittelbar das Fundamentaltheorem folgt.

\clearpage\SourceLeaf{174}%
Das Fundamentaltheorem kann zusammen mit einigen bekannten Sätzen zu einer ziemlich kurzen Lösung der Aufgabe dienen, zu entscheiden, ob eine gegebene positive ganze Zahl quadratischer Rest oder Nichtrest einer gegebenen Primzahl ist, wie in der Abhandlung ausführlich gezeigt wird. Weiteres Nachdenken über den dritten Beweis des Fundamentaltheorems führte den Verfasser jedoch zu einer noch viel geschmeidigeren Lösung. Sie bildet den dritten Teil der Abhandlung. Hier geben wir nur ihre abschließende Regel an und lassen die Herleitung aus Gründen der Kürze weg. Soll entschieden werden, ob die positive ganze Zahl $b$, die durch die Primzahl $a$ nicht teilbar ist, quadratischer Rest oder Nichtrest von dieser ist, so bilde man genau wie bei der Suche nach dem größten gemeinsamen Teiler von $a$ und $b$ die Gleichungen
\[
\begin{aligned}
a&=\beta b+c\\
b&=\gamma c+d\\
c&=\delta d+e\\
d&=\varepsilon e+f\text{ usw.}
\end{aligned}
\]
bis man in der Folge der Zahlen $a,b,c,d,e,f$ usw. auf .die Einheit kommt. Man. bezeichne die Zahlen $\frac12a,\frac12b,\frac12c,\frac12d$ usw., unter Weglassung ihres Bruchteils $\frac12$, sofern einige der Zahlen $a,b,c,d$ usw. ungerade sind, mit $a',b',c',d'$ usw. Man nenne $\mu$ die Anzahl der in der Folge $a',b',c',d'$ usw. vorkommenden Paare unmittelbar aufeinanderfolgender ungerader Zahlen. Schließlich nenne man $\nu$ die Anzahl der ungeraden Zahlen in der Folge $\beta,\gamma,\delta,\varepsilon$ usw., denen in der Folge $b',c',d',e'$ usw. der Reihe nach eine Zahl der Form $4k+1$ oder $4k+2$ entspricht. Dann ist $b$ quadratischer Rest oder Nichtrest von $a$, je nachdem $\mu+\nu$ gerade oder ungerade ist, mit der einzigen Ausnahme, dass zugleich $b$ gerade ist und $a$ die Form $8k+3$ oder $8k+5$ hat. In diesem Fall gilt die entgegengesetzte Regel: Ein gerades $\mu+\nu$ zeigt an, dass $b$ quadratischer Nichtrest von $a$ ist, ein ungerades $\mu+\nu$ dagegen, dass $b$ quadratischer Rest von $a$ ist.
\BookRule{.28\linewidth}
